To display a graph you can bind a graphic context to a graph.
Graphic contexts are C++ classes which are subclasses of {\bf
  FFGraphWindow} which is a subclass of {\bf FFGraphLabel}. How
the graph is displayed depends on the graphic context you use.
This may range from a window on the screen to an output to a
printer. When a graphic context is bound to a graph, the graph is
automatically labeled with the node label {\bf FFGraphicNode} and
the edge label {\bf FFGraphicEdge} which contain graphic
information like coordinates, colors, etc. If you just want the
graphic information but do not want to actually display the graph,
e.g. to do some layouting, you can use the special graphic
context {\bf NULL}.

The library contains the class {\bf FFTclWindow} which is able to
display 2 and 3 dimensional graphs. 

\begin{figure}[htb]
\fbox{\epsfig{figure=icons.eps,clip=}}
\caption{Graph containing icons}
\end{figure}

It is also able to display nested graphs:

\clearpage

\begin{figure}[H]
\hspace*{-3cm}
\fbox{\epsfig{figure=g5.eps,clip=}}
\caption{Graph containing a subgraph}
\end{figure}
\clearpage
\subsection{Using graphic contexts} 
\label{using_graphic}


To assign a graphic context to a graph you should use the method \\

\begin{itemize}
\item{\bf void FFGraphClass::createGraphic(FFGraphWindow *window)}. \\
\end{itemize}

E.g. to display a graph using the graphic context {\bf FFTclWindow} you 
could write: 

\bigskip

\begin{verbatim}
FFGraph g;
g->createGraphic(new FFTclWindow);
\end{verbatim}

\bigskip

If you want to use graphic information without actually
displaying the graph, you can call createGraphic with {\bf window
  == NULL}.

If a graphic context is assigned to a graph, it can be retrieved
using the method\\

\begin{itemize}
\item{\bf FFGraphWindow *FFGraphClass::getWindow()}\\
\end{itemize}

At the moment only one subclass of {\bf FFGraphWindow} can be assigned
to a graph. If you no longer want to display the graph, you can use the
method\\

\begin{itemize}
\item{\bf void FFGraphClass::removeGraphic()}\\
\end{itemize}

which removes the graph-window label and the labels {\bf FFGraphicNode}
and {\bf FFGraphicEdge}. The labels are freed using the {\bf delete}
operator. If you do not want the graph window label to be freed, e.g. to
use it for another graph, you can use the method\\

\begin{itemize}
\item{\bf FFGraphWindow *FFGraphClass::leaveWindow()}\\
\end{itemize}

\subsection{Writing graphic contexts}
\label{writing_graphic}

To write your own graphic context, you can either write a new one from
scratch, or subclass an existing one and only overwrite the methods you
want to change. Graphic contexts are designed to be able to display 3
dimensional graphs. If you do not want to display 3D graphs, just
implement the appropriate methods to ignore 3D information.  A graphic
context must be a subclass of {\bf FFGraphWindow} and must implement all
methods {\bf FFGraphWindow} defines.  See section \vref{ffgraphwindow}
for details on what methods you have to implement. When {\bf
  FFGraphClass::createGraphic(...)} is invoked with your window label as
argument, the graph labels itself with the node label {\bf
  FFGraphicNode} and the edge label {\bf FFGraphicEdge}. You should use
these labels to store and get the graphic information you need to
display the graph. See sections \vref{ffgraphicnode} and
\vref{ffgraphicedge} for further information on how to use these labels.

The easiest way to write your own graphic context is to subclass {\bf
FFTclWindow}. Of course by doing so you also decide to use the Tcl/Tk
widget library. As this library is very ({\em VERY}) slow it might
be the easiest way but not necessarily the best. Anyway it is a good
starting point and I use Tcl/Tk myself. 

Ok, let's assume you decided to implement your graphic context based
on {\bf FFTclWindow}. Of course you can now overload each method and
write a completely new tcl script, but most probably 90\% of the
window are suitable for your purposes anyway. Usually you only want to
change the events, add a menu, change some controls, etc. Therefore I
added some support for doing exactly this. This part will tell you
how to buid a class that behaves exaclty like {\em FFTclWindow}. Using
this class as basis you then will learn what possibilities you have to
change the look and behavior. For many things a simple mechanism is
provided. If you intend to do something which is not within the scope
of the existing mechanism you might have to rewrite bigger parts of
the standard tcl script. You can also drop me a note and if I think
it's a good idea it will probably be a feature in the next version.

As I noted before the Tcl graphic contexts merely are Tcl interpreters
which execute a Tcl script which implements the functionality of the
window. Consequently there are basicly two points where you can change
the behavior of a Tcl graphic context. The first is the C++ source
code of your subclass and the second is the Tcl-script for your class.
 
The simplest possible graphic context might look like this:

{\tt

class MyWindow: public FFTclWindow
\{
public:
 virtual FFGraphLabel *proto() const \{return new MyWindow;\}
\};
}

This context will look and behave exactly like FFTclWindow. You now
can overload methods and so change the behavior of the Window. As the
most of the graphic part is defined in the Tcl-script you won't be
able to change much of it in the C++ file. So what you have to do is to tell
your class to use your onw Tcl-script and of course you have to write
one. As you can reuse much of the existing Tcl-stuff this is not much
work as well. 

First you have to set the Membervariable {\em tclFileName} to point to
your Tcl-script. A good place to do this is the constructor of your
class. If you want to use your own icons, you should set {\em
iconLibDir} as well:

\begin{verbatim}
MyWindow::MyWindow()
{

 // Set the Scriptname
 
 String tmpS;
 
 if(getenv("FF_LIBDIR") == NULL)
  tmpS=String(FF_LIBDIR);
 else
  tmpS=String(getenv("FF_LIBDIR"));
 
 tclFileName =tmpS + "/tcl/myscript.tcl";

 // Set the icon location     

 if(getenv("FF_ICONDIR") == NULL)
  tmpS=String(FF_LIBDIR) + "/icons/";
 else
  tmpS=String(getenv("FF_ICONDIR"));

 iconLibDir = tmpS;
}
\end{verbatim}

The script itself can be as simple as:

\begin{verbatim}
#!/public/bin/wish -f

global FF_LIBDIR
source $FF_LIBDIR/tcl/xfmisc.tcl
source $FF_LIBDIR/tcl/fileselect.tcl
source $FF_LIBDIR/tcl/stdGraph.tcl

ShowWindow.
\end{verbatim}


Again you have a class which does look and behave exactly like
FFTclWindow. We now have a Tcl script and a C++ class, but how does
communication take place? 

Tcl functions can be called from C++ using the C functions descibed in
the Tcl/Tk manual like {\em Tcl_VarEval} etc. The interpreter you have
to supply is stored in the member variable {\em interp}.

If the tcl command {\em recall} is called the method
{FFTclBasicWindow::recall(Tcl_Interp *i, int c, char **v)} is
invoked. {\em recall} can have an arbitrary number of arguments. {\em
c} is the number of arguments and {\em v} is a vector containing the
arguments. If you want to add your own commands or change the way
existing commands are handled you can overload the method {\en
recall}:

\begin{verbatim}
int MyWindow::recall(Tcl_Interp *i, int c, char **v)
{
 String command(v[1]); // v[0] == ``recall''

 if( command == myNewCommand)
  {
    doSomething();
    return TCL_OK;
  }
 else if (command == commandIWantToHandleDifferent)
  {
    doSomethingElse();
    return TCL_OK;
  }
 else
  return FFTclWindow::recall(i,c,v);

 return TCL_OK;
}
\end{verbatim}

You now know what you can do on the C++ side to change the
behavior. The next part will deal with the question what you can do on
the Tcl side.

A graph window as defined in {\em stdGraph.tcl} consists of four
parts. The first is the main menu which is at the top of the window.
Then comes a toolbar which is not visible in {\em FFTclWindow}. Next
is the area where the rotation and scalation scales are located. At
the bottom of the window is the graph canvas. Each of these parts can
be modified.

\begin{itemize}

\item Changing the main menu.

If you want to add menu entries you can do so by overloading the tcl
function {\em addUserMenus \{frameArg\}}. {\em frameArg} contains the
name of the menubar frame. In this function you can create new
menubuttons and append them to the frame.

There is currently no way to avoid menu entries you don't want. The
only way to do this is to rewrite the Tcl function {\em
buildMainMenu}. This function is expected to create the frame {\em
.frame1} and build the menu in this frame.

\item Changing the Toolbar

The toolbar can be modified by rewriting the Tcl function {\em
buildToolBar}. This function is expected to create the frame {\em
.toolFrame} and build whatever it wants inside this frame.

\item Changing the usage box

The usage box is shown if the user activates the {\em help} menu. The
help box is a toplevel window. The window is passed to the Tcl function
{\em addUserUsage\{windowName\}}. You can overwrite this function and
put you help information in this window.

\item Changing the scale area

The scale area can be modified by rewriting the Tcl function {\em 
buildButtonFrame}. This function is expected to create the frame {\em
.frame2}. The function is free to build whatever it wants inside.

\item Changing the canvas area

You are not intened to change the look of this area. At least in the
current release. What you can easily change is the way the canvas
reacts on certain events. Events are interactions by the user like
pressing a key, moving the mouse etc. There are three different kinds
of events. Events can be bound to the canvas, to nodes and to edges. 
Usually events are associated with the canvas. If a node or an edge is
selected the event is associated with the node or edge. 

You can overload the functions {\em addUserLineBindings}, {\em
addUserNodeBindings} and {\em addUserCanvasBindings} to add events you
want to react on. You can delete an event by redfining it to do
nothing. 

\end{itemize}
